\section{Avenues for further research}
\label{sec:future}

Each avenue below addresses the limitation carrying the same number in
Section~\ref{sec:limitations}.

\begin{description}[leftmargin=2.4em,style=nextline]

  \item[FR1] Researchers should source the three swept components directly, the rent share from
  firm-level production data, the deferral factor from realization histories and the debt share
  from a marginal rather than a stock ratio, which would replace the swept grid with an estimate
  carrying a sampling distribution and turn our share of a space into a confidence statement.

  \item[FR2] Researchers should construct both bases for the taxable shareholder share from the
  same underlying holder data and report the assembled rate on each, rather than taking one base
  from one literature and comparing it with a figure built on the other, which is the comparison
  our own published check had to be corrected for.

  \item[FR3] Researchers should measure the effective withholding rate actually borne by foreign
  holders treaty by treaty, and separate dividends, which bear withholding, from capital gains
  realized by foreign holders, which generally do not. That would replace the single stated
  treaty rate with a measured one and close the largest open assumption in the base construction.

  \item[FR4] Researchers should build an annual-revenue counterpart to the marginal wedge, using
  an investment path and a capital stock rather than a single marginal project, so the quantity
  compared against a replacement requirement is the quantity a budget actually collects. Our
  illustrative path shows the mechanism and does not size it.

  \item[FR5] Researchers should replace the single economy-wide labour tax rate with rates by
  position in the income distribution, and pair them with an assumed incidence of displacement
  across that distribution, which would sign and bound an effect we can currently only describe.

  \item[FR6] Researchers should estimate the pass-through coefficient directly, measuring how
  much taxable domestic capital income a dollar of displaced wages actually generates, which is
  the single input that would convert our boundary from a surface into a located point. We
  located no estimate of it anywhere, which may itself be the finding.

  \item[FR7] Researchers should embed the incidence arithmetic in a model that lets households
  adjust, so the measured boundary can be compared against a behavioural one and the difference
  attributed. \citet{Bayraktar2026} supplies one such model and the comparison has not been made.

  \item[FR8] Researchers should resolve the housing agency classification against the legal and
  budgetary treatment of conservatorship rather than reporting both readings, which would remove
  the factor-of-seven swing that is currently the largest unresolved judgement in the incidence
  result.

  \item[FR9] Researchers should extend the bust bracket beyond two US episodes, using
  international equity busts with comparable fiscal detail, and should estimate the wealth effect
  on the relevant population rather than importing a sourced average.

  \item[FR10] Researchers should build the off-balance-sheet side of the capital stock from
  securitization prospectuses, vehicle-level filings and supervisory data on lending secured
  against computing hardware, which would make the financing verdict unconditional rather than
  restricted to what nine balance sheets disclose.

\end{description}
